\documentclass[11pt]{article}
\usepackage[a4paper,margin=25mm]{geometry}
\usepackage{amsmath,amssymb,amsthm}
\usepackage[hidelinks]{hyperref}
\newtheorem{theorem}{Theorem}
\theoremstyle{remark}
\newtheorem*{remark}{Remark}
\title{Two Eight-Element Distributive Lattices:\\A Disproof of Conjecture 00000002604}
\author{GitHub: gaochengzhecpu}
\date{}
\begin{document}
\maketitle
\begin{abstract}
The source asserts that finite distributive lattices with the same
layer-count vector are isomorphic. The lattices $C_4\times C_2$ and
$(C_3\times C_3)\setminus\{(2,0)\}$ are distributive, have eight
elements each, and have layer-count vector $(1,2,2,2,1)$ whether layers
are counted from the bottom or from the top. They are not isomorphic.
All of this, and the negation of the universal claim, is verified in
Lean. An earlier submission used an equivalent example but did not
establish non-isomorphism; its error is described below.
\end{abstract}

\paragraph{Reading of the source.}
A finite distributive lattice is graded: it has a unique rank function
$\rho$ with $\rho(\bot)=0$ and $\rho(b)=\rho(a)+1$ whenever $b$ covers
$a$. Its $k$-th layer is $\{x:\rho(x)=k\}$. The English text speaks of
the layer-count vector; the Chinese text counts the layers from the
coatoms downwards, that is by the corank $\rho(\top)-\rho(x)$. We treat
both. Layer $1$ of the corank is the set of coatoms.

\begin{theorem}
Let $C_n$ be the chain $0<1<\dots<n-1$, and let
\[
 L_1=C_4\times C_2,\qquad
 L_2=\{(i,j)\in C_3\times C_3:\ (i,j)\ne(2,0)\},
\]
both with the componentwise order. Then $L_1$ and $L_2$ are distributive
lattices with eight elements. Their layer-count vectors, by rank and by
corank, are all equal to $(1,2,2,2,1)$. The lattices $L_1$ and $L_2$ are
not isomorphic.
\end{theorem}
\begin{proof}
\emph{Lattices.} A product of chains is a distributive lattice under
componentwise maximum and minimum. Let $p,q\in L_2$. If $p\vee q$ has
first coordinate $2$, then $p$ or $q$ has first coordinate $2$ and so a
nonzero second coordinate; hence $p\vee q$ has a nonzero second
coordinate. If $p\wedge q$ has first coordinate $2$, then both $p$ and
$q$ do, both second coordinates are nonzero, and so is their minimum.
So $L_2$ is a sublattice of $C_3\times C_3$ and is distributive.

\emph{Layers.} In both lattices $\rho(i,j)=i+j$. This is clear for
$L_1$. For $L_2$, every pair $a<b$ of $L_2$ that has $(2,0)$ strictly
between also has $(1,0)$ or $(1,1)$ strictly between. Removing $(2,0)$
therefore creates no new covering relation, and every covering relation
of $L_2$ raises $i+j$ by one. The layers are
\[
\begin{array}{c|ccccc}
 i+j&0&1&2&3&4\\\hline
 L_1&(0,0)&(1,0),(0,1)&(2,0),(1,1)&(3,0),(2,1)&(3,1)\\
 L_2&(0,0)&(1,0),(0,1)&(1,1),(0,2)&(2,1),(1,2)&(2,2)
\end{array}
\]
so both vectors are $(1,2,2,2,1)$. This vector is a palindrome and the
top element has rank $4$ in both lattices, so the corank layers give the
same vector. In particular both lattices have two coatoms.

\emph{Non-isomorphism.} Call an element a \emph{branch point} if two
incomparable elements lie below it and two incomparable elements lie
above it. This property is preserved by order isomorphisms. In $L_2$ the
element $(1,1)$ is a branch point: $(1,0)$ and $(0,1)$ lie below it, and
$(2,1)$ and $(1,2)$ lie above it. In $L_1$ the elements below $(i,0)$
form a chain, and the elements above $(i,1)$ form a chain. So $L_1$ has
no branch point, and no isomorphism $L_1\to L_2$ exists.
\end{proof}

\begin{remark}
By Birkhoff's theorem, $L_1$ is the lattice of order ideals of the
poset $P$ formed by a three-element chain and an isolated point, and
$L_2$ is that of the poset $Q$ on $\{a,b,c,d\}$ with $a<c$, $b<c$,
$b<d$. For $Q$ the isomorphism sends an ideal $I$ to
$(|I\cap\{a,c\}|,\,|I\cap\{b,d\}|)$; the excluded pair $(2,0)$ would be
an ideal containing $c$ but not $b$. This remark is not formalised and
is not needed.
\end{remark}

\section*{The earlier submission and its error}
Pull request 48 of the competition repository (by \texttt{earthking11},
merged on 1 October 2026) treated this conjecture with the posets $P$
and $Q$ of the remark. The organisers removed it in the re-audit commit
\texttt{541cf4f} of 3 October 2026, with the note that non-isomorphism
was never established.

Its mathematics is sound: the two ideal lattices do have equal vectors
and are not isomorphic. The error is in the Lean file. Its final theorem
asserts
\begin{center}\small
\texttt{perms.all (fun p => IsIsoPerm leP leQ p) = false}.
\end{center}
This says that it is false that all $24$ bijections are isomorphisms,
in other words that at least one bijection is not an isomorphism. It
does not say that no bijection is an isomorphism, which is what
non-isomorphism means. In addition, the file works with Boolean
relations on four points and with bitmask counts of ideals. It never
forms the two lattices, their order, or a statement about lattice
isomorphism.

The present submission forms both lattices as Mathlib distributive
lattices and proves that the type of order isomorphisms between them is
empty.

\section*{Formal verification and scope}
\texttt{L1} is the product type of \texttt{Fin 4} and \texttt{Fin 2}
with Mathlib's product order, which carries Mathlib's
\texttt{DistribLattice} instance. \texttt{L2} is the subtype of the
product of \texttt{Fin 3} with itself defined by the condition above. Closure under supremum and infimum is proved, and the
\texttt{DistribLattice} instance of the subtype is derived from it, so
its operations are those of the ambient lattice. Both carry bounded
orders.

\texttt{IsRankFunction} and \texttt{IsCorankFunction} are stated with
Mathlib's covering relation. The functions $i+j$ and $4-(i+j)$ are
proved to satisfy them on both lattices. The eight equalities of layer
counts are kernel-checked, and \texttt{layer\_rank\_eq} and
\texttt{layer\_corank\_eq} give equality of the layer counts for every
$k$. \texttt{not\_isomorphic} proves that the type of order isomorphisms
from \texttt{L1} to \texttt{L2} is empty, by the branch-point argument.

The proposition
\begin{center}\small\texttt{ClaimedCoatomLayerRigidity}\end{center}
quantifies over all pairs of finite bounded distributive lattices with
corank functions and equal layer counts and asserts an order
isomorphism. Its negation is \texttt{conjecture\_false}. The same
statement for rank functions is refuted in
\texttt{rank\_rigidity\_false}. No \texttt{sorry},
\texttt{native\_decide} or custom axiom occurs.

The disproof concerns the first clause of the source. Since two
non-isomorphic lattices share a vector, no characterisation of the
realisable vectors can repair that clause, and no claim is made about
the Kruskal--Katona clause.

\section*{Source and reproduction}
The exact bilingual source is included byte-for-byte as
\texttt{SOURCE.md}. Its repository path is
\begin{center}\small\texttt{conjectures/00000002604.md}.\end{center}
The project pins Lean 4.19.0 and Mathlib commit
\begin{center}\small\texttt{c44e0c8ee63ca166450922a373c7409c5d26b00b}.\end{center}
From \texttt{lean/}, run \texttt{lake build}, then
\begin{center}\small\texttt{lake env lean -DwarningAsError=true Main.lean}.\end{center}
The included public-Git manifest locks all dependency revisions. The
\texttt{verification/} directory records the fresh-build logs,
theorem-axiom reports, artifact hashes, review and final PDF
inspection. No supplementary numerical computation is required.
\end{document}
